\begin{proof}
Let  

\[
\mathbf a=(a_{1},\dots ,a_{n}),\qquad   
\mathbf b=(b_{1},\dots ,b_{n})\in\mathbb R^{n}.
\]

Define the Euclidean inner product and norm  

\[
\langle\mathbf a,\mathbf b\rangle =\sum_{i=1}^{n}a_{i}b_{i},\qquad 
\|\mathbf a\|^{2}= \langle\mathbf a,\mathbf a\rangle =\sum_{i=1}^{n}a_{i}^{2},\qquad 
\|\mathbf b\|^{2}= \langle\mathbf b,\mathbf b\rangle =\sum_{i=1}^{n}b_{i}^{2}.
\]

\medskip\noindent\textbf{Degenerate cases.}  
If \(\mathbf a=\mathbf 0\) or \(\mathbf b=\mathbf 0\) then  
\(\langle\mathbf a,\mathbf b\rangle =0\) and also \(\|\mathbf a\|^{2}=0\) or \(\|\mathbf b\|^{2}=0\). Hence  

\[
\bigl(\langle\mathbf a,\mathbf b\rangle\bigr)^{2}=0=
\|\mathbf a\|^{2}\,\|\mathbf b\|^{2},
\]

so the inequality holds with equality.  
Henceforth assume \(\mathbf a\neq\mathbf 0\) and \(\mathbf b\neq\mathbf 0\); consequently  
\(\|\mathbf a\|^{2}>0\) and \(\|\mathbf b\|^{2}>0\).

\medskip\noindent\textbf{A non‑negative quadratic.}  
For a real parameter \(t\) define  

\[
Q(t)=\sum_{i=1}^{n}\bigl(a_{i}-t b_{i}\bigr)^{2}.
\]

Each summand is a square, therefore \(Q(t)\ge 0\) for every \(t\in\mathbb R\). Expanding,

\[
\begin{aligned}
Q(t)&=\sum_{i=1}^{n}\bigl(a_{i}^{2}-2t a_{i}b_{i}+t^{2}b_{i}^{2}\bigr)\\[2mm]
    &=\|\mathbf a\|^{2}-2t\langle\mathbf a,\mathbf b\rangle
      +t^{2}\|\mathbf b\|^{2}.
\end{aligned}
\]

Thus  

\[
Q(t)=\|\mathbf b\|^{2}\,t^{2}-2\langle\mathbf a,\mathbf b\rangle\,t+\|\mathbf a\|^{2},
\]

a quadratic polynomial with positive leading coefficient \(\|\mathbf b\|^{2}>0\).

\medskip\noindent\textbf{Discriminant condition.}  
A quadratic \(At^{2}+Bt+C\) with \(A>0\) satisfies \(At^{2}+Bt+C\ge0\) for all real \(t\) iff its discriminant is non‑positive:

\[
B^{2}-4AC\le0 .
\]

Applying this to the coefficients  

\[
A=\|\mathbf b\|^{2},\qquad B=-2\langle\mathbf a,\mathbf b\rangle,\qquad 
C=\|\mathbf a\|^{2},
\]

we obtain  

\[
\Delta =(-2\langle\mathbf a,\mathbf b\rangle)^{2}
        -4\|\mathbf b\|^{2}\|\mathbf a\|^{2}
      \le0 .
\]

Since  

\[
(-2\langle\mathbf a,\mathbf b\rangle)^{2}
-4\|\mathbf b\|^{2}\|\mathbf a\|^{2}
=4\bigl(\langle\mathbf a,\mathbf b\rangle^{2}
        -\|\mathbf a\|^{2}\|\mathbf b\|^{2}\bigr),
\]

dividing the inequality \(\Delta\le0\) by the positive constant \(4\) yields  

\[
\langle\mathbf a,\mathbf b\rangle^{2}\le\|\mathbf a\|^{2}\,\|\mathbf b\|^{2}.
\]

Writing the inner products and norms in component form gives the **Cauchy–Schwarz inequality**

\[
\boxed{\displaystyle
\biggl(\sum_{i=1}^{n}a_{i}b_{i}\biggr)^{2}
\le
\biggl(\sum_{i=1}^{n}a_{i}^{2}\biggr)
   \biggl(\sum_{i=1}^{n}b_{i}^{2}\biggr)}.
\]

\medskip\noindent\textbf{Equality condition.}  
Equality holds precisely when \(\Delta=0\), i.e. when the quadratic has a double root  
\(t_{0}= \dfrac{\langle\mathbf a,\mathbf b\rangle}{\|\mathbf b\|^{2}}\).  
Since \(Q(t_{0})=0\) and each term \((a_{i}-t_{0}b_{i})^{2}\ge0\), we must have  

\[
a_{i}=t_{0}b_{i}\qquad (i=1,\dots ,n).
\]

Thus the vectors are linearly dependent; there exists \(\lambda\in\mathbb R\) such that  
\(a_{i}= \lambda b_{i}\) for every \(i\) (the trivial case \(\lambda=0\) corresponds to one of the vectors being the zero vector).

Hence equality in the Cauchy–Schwarz inequality holds **iff** the two sequences are proportional.

\end{proof}